\section{How many products can be prescribed at once?}\label{sec:growing}
The article asks how fast the number of factors can grow at retained density at most one half. For one common modulus and an unrestricted residue in the reserve block, the answer has the same order as the elementary capacity bound. The result in this section allows every fixed density \(0<a\le1\). All constants are independent of the number of target blocks.

\begin{theorem}[Growing number of targets at any positive density]\label{th:growing}
Fix \(0<a\le1\) and \(A>0\). There are \(C=C(a)\), \(B=B(a,A)\) such that the following holds under Dirichlet GRH for every sufficiently large \(q\). Take any \(S\subseteq\Pp\) of size \(m\ge aN\) and any \(r\ge1\) for which the quotas, including the reserve, satisfy \(h\ge B\log q\). Put \(D=\lfloor1/a\rfloor\). There is a group \(H\subseteq\widehat G\) of size \(d\le D\), depending only on \(q,S\), with evaluation quotient \(\pi:G\to\widehat H\). Write \(\nu\) for the joint law of the \(r\) projected products. The support is exactly the inverse image of \(\operatorname{supp}\nu\), and at each occupied target,
\[
  \pr(X=g)=\frac{d^r}{|G|^r}\nu(\pi^r(g))(1+\theta_g),
       \qquad |\theta_g|\le q^{-A}.
\]
For \(a>1/4\), the target tests of Theorem~\ref{th:support} and the coverage thresholds \eqref{eq:coverage} hold throughout this growing range. Balanced quotas allow \(r\asymp_{a,A} \log q/\log\log q\) even at the least allowed retained density. The number of primes available makes this order maximal under the stipulated quota size.
\end{theorem}

\begin{proof}
The Gram argument after \eqref{eq:pool} works verbatim with any fixed \(D\). Choose \(\tau=\tau(a)\), then \(C\), so the subgroup generated by biases exceeding \(1-\tau\) has size \(d\le D\). Put \(k=r+1\). For \(B\ge1\), \(k\le m/\log q\). From \eqref{eq:pool} and \(m\le N\) there is a fixed \(K=3(C+1)\) such that
\begin{equation}\label{eq:growsize}
 k\log(m+1)\le K\log q .
\end{equation}
Indeed \(N\sim C(\log q)^2/\log(C(\log q)^2)\). We give two estimates whose losses remain bounded by fixed powers of \(q\).

\emph{The Fourier error.}
In the Cauchy estimate leading to \eqref{eq:lift}, retain every pair of labels whose characters lie in different cosets of \(H\). Give each coset the weight \(w_c=\sum_{i:\psi_i\in c}n_i/m\), including \(\psi_0=1\). With \(c_*\) a coset of largest weight and \(t_*\) the number of labels outside it,
\[
 \sum_{c<c'}w_cw_{c'}=\tfrac12(1-\textstyle\sum_c w_c^2)
 \ge \tfrac12(1-w_{c_*})\ge t_* h/(2m).
\]
Thus the Fourier coefficient is bounded by \((m+1)^k e^{-\tau t_* h/2}\). A tuple outside \(H^r\) has \(t_*\ge1\), since the reserve character is trivial. There are at most \(q\binom{k}{t_*}d^k q^{t_*}\) tuples for a given \(t_*\): choose the dominant coset, the other labels, and their characters. This count permits every possible reserve character and hence majorizes the required count. If \(B\ge8/\tau\), then \(kq e^{-\tau h/2}\le q^{-2}\) for large \(q\). Summing the binomial bound, using \((1+x)^k-1\le2kx\) for \(kx\le1/2\), bounds the total error by
\begin{equation}\label{eq:growerror}
 (m+1)^k q d^k\big[(1+q e^{-\tau h/2})^k-1\big]
 \le 2q^{2K+3-\tau B/2}.
\end{equation}
Here \(k\le q\) and \(d\le m+1\) for large \(q\). Orthogonality on \(H^r\) gives the principal term \(d^r\nu(\pi^r(g))\) in the normalization of \eqref{eq:lift}.

\emph{Repair of the color counts.}
Call \(\pi(p)\in\widehat H\) the color of a prime, and let \(e\le d\) be the exponent of \(\widehat H\). Fix one allocation witnessing an occupied projected tuple. Its \(k\)-by-\(d\) table \(Y\) records the number of each color in each block. The table \(X'\) of any other allocation has the same row and column sums. We construct a table \(X''\) with those margins and with every entry congruent to the corresponding entry of \(Y\) modulo \(e\). This ensures that each block has exactly the witnessed color product.

Let \(\sigma_{ij}\in\{0,\ldots,e-1\}\) be the residue of \(Y_{ij}\). Take the row and column sums of \((Y-\sigma)/e\) as target margins. Start at the nonnegative integer table
\[
 Z_{ij}=\max\{0,\lfloor(X'_{ij}-\sigma_{ij})/e\rfloor\}.
\]
Initially \(\|\sigma+eZ-X'\|_1\le e k d\), where the norm sums absolute entries. Both the total absolute row discrepancy and the total absolute column discrepancy of \(Z\) from the target margins are at most \(kd\), by the equal margins of \(X',Y\). Delete units from overfull rows until none remains, and then from overfull columns until none remains. Each stage costs at most \(kd\), as row deletions can only decrease column excesses. The remaining summed row deficit equals the summed column deficit and is at most \(3kd\). Fill it by adding units at intersections of deficient rows and columns. Writing \(Z_*\) for the repaired table, we obtain
\[
 X''=\sigma+e Z_*,\qquad \|X''-X'\|_1\le 6e k d\le6D^2 k.
\]
For each color, transfer individual primes from its surplus blocks to its deficient blocks. This changes at most \(6D^2 k\) prime labels, preserves every quota in the final allocation, and reaches the witnessed projection. The number of allocations within that many changes of any one allocation is at most \((1+mk)^{6D^2k}\): pad a list of edits with empty moves. Since all allocations have the same probability, the neighborhoods of the allocations realizing the projection cover the full space. Thus
\[
  \nu(\pi^r(g))\ge (1+mk)^{-6D^2 k}\ge q^{-12D^2 K}
\]
at every occupied point by \eqref{eq:growsize}. Take, for instance,
\(B\ge (2/\tau)(A+12D^2K+2K+6)\), also with \(B\ge1\). Dividing \eqref{eq:growerror} by the occupied main term gives the stated relative error. An empty projection rules out every lift by definition.

Finally, \eqref{eq:pool} and \(kh\le m\) imply \(r=O_{a,A}(\log q/\log\log q)\), hence \(h\ge2r\) uniformly for large \(q\). When \(a>1/4\) the quotient has order at most three. The finite sign and cubic allocation arguments of Theorem~\ref{th:support} apply unchanged: the quotas can hold \(2r\) exceptional primes and a most frequent class supplies at least \(m/3\ge2r\) anchors. Balanced quotas with \(k=\lfloor aN/(2B\log q)\rfloor\) demonstrate the claimed attainable order for every permitted retained set.
\end{proof}

The quotient in this theorem uses a single unit group \(G\). The next section allows an independently chosen modulus for each block.
